\section*{Asesmen}
\addcontentsline{toc}{section}{Asesmen}

Bab VII memuat Praktikum 1 dengan bobot 12\% sesuai RPS. Asesmen
mencakup kebenaran, kualitas kode, kelengkapan, dan pengujian.

\begin{table}[H]
\centering
\caption{Instrumen asesmen Bab VII}
\label{tab:ch07-asesmen}
\begin{tabularx}{\textwidth}{@{}lXl@{}}
\toprule
\textbf{Instrumen} & \textbf{Indikator} & \textbf{Bobot} \\
\midrule
Praktikum 1 & Lexer tangan dan Flex, 30 tes, keluaran uji, laporan &
12\% \\
Kuis & Buffer, maximal munch, pemulihan & formatif \\
Review kode & Keterbacaan, struktur, penanganan kesalahan & bagian P1 \\
\bottomrule
\end{tabularx}
\end{table}

Kriteria kelulusan: lexer mengenali seluruh token Subset-C, lulus
minimal 30 tes, menghasilkan diagnostik yang informatif, dan
terdokumentasi dengan baik.
